\section{Limitaciones y conclusión}
\label{sec:conclusion}

\paragraph{Limitaciones.}
Los dos bancos son franceses, así que la réplica varía año, región, reparto de clases y familia de
modelo, pero no el sistema agrícola; queda abierto si la descomposición se sostiene en un paisaje
de agricultura de pequeña escala. En el banco de réplica los bloques son las dos regiones
administrativas, porque el índice publicado no trae coordenadas por parcela, así que nuestro
remuestreo captura la variación dentro de las regiones y no entre ellas; por eso reportamos las
diferencias por bloque. Ese banco exigió además imputar el \num{1.62}\,\% de las celdas de
características, concentradas en índices de cociente cuyo denominador involucra una banda que la
colección codifica como cero para el píxel enmascarado; la imputación usa solo la mediana de la
región de entrenamiento, y no puede favorecer a un mecanismo sobre otro porque todos leen la misma
posterior. Por último, nuestras posteriores fuera de fold en el banco primario existen para un solo
fold, y por eso el meta-modelo de la sección~\ref{sec:arbitro} se entrena con pocos datos; eso es
una limitación de los artefactos de que disponemos, no una afirmación sobre el apilamiento en
general.

\paragraph{Conclusión.}
El F1 promediado por clase sobre un catálogo podado mezcla dos cosas que quien lee no puede
separar: lo que hizo el método y lo que hizo el denominador más pequeño. Dimos un protocolo que las
separa, y la respuesta en dos bancos públicos es que manda el denominador: el \num{88}\,\% en el
primero, el \num{95}\,\% en el segundo en el punto de operación preregistrado, y el \num{100}\,\%
en el tramo donde la cobertura sigue completa y todos los mecanismos devuelven una cifra idéntica.
A igual cobertura no encontramos evidencia de que recortar la leyenda supere a abstenerse, y sí
evidencia significativa de lo contrario en el banco más desbalanceado; el criterio que se reporta
usar, la retirada por soporte bajo, es el más débil que probamos y empeora conforme se retiran más
clases. Nada de esto exige un modelo nuevo. Exige reportar el número de clases, la cobertura y una
línea de control junto a cada cifra macro calculada sobre un catálogo que alguien eligió.
